\documentclass[10pt]{letter}

\usepackage[T1]{fontenc}
\usepackage[utf8]{inputenc}
\usepackage[margin=0.75in]{geometry}
\usepackage{microtype}
\usepackage[hidelinks]{hyperref}

\pagestyle{empty}
\setlength{\parskip}{0.45em}
\setlength{\longindentation}{0pt}

\address{%
Jeffrey D. Varner\\
Robert Frederick Smith School of Chemical and Biomolecular Engineering\\
Cornell University\\
Ithaca, New York 14850, United States\\
\href{mailto:jdv27@cornell.edu}{jdv27@cornell.edu}}
\signature{Jeffrey D. Varner}
\date{August 11, 2026}

\begin{document}

\begin{letter}{%
Editors\\
\textit{Journal of Chemical Information and Modeling}}

\opening{Dear Editors:}

Please consider our Article, ``Conditioning Protein Generation via Hopfield Pattern
Multiplicity,'' for publication in the \textit{Journal of Chemical Information and
Modeling}.

This work presents a training-free method for conditioning protein sequence generation when
only a small designated subset of a family is available. Pattern multiplicities enter the
modern Hopfield energy as logit biases, giving an exact Gaussian-mixture equilibrium target
without retraining or changing the sampler's asymptotic cost. The paper also distinguishes
this latent target from the sequence-level distribution obtained after PCA reconstruction and
discrete decoding, exposing and quantifying a practically important calibration gap.

The computational study includes five Pfam families and a curated $\omega$-conotoxin family,
replicated multiplicity sweeps, exact-equilibrium controls, hard-mask and hard-curation
comparisons, and sequence, structure, and language-model assessments. We additionally
introduce a matched HMMER3 benchmark that uses the same cleaned alignments, designation
labels, multiplicity ratios, replicate counts, and output sizes. The comparison gives an
informative tradeoff rather than a blanket superiority claim: the weighted profile HMM more
directly transfers a selected position-specific marginal, whereas stochastic attention gives
lower ESM2 pseudo-perplexity in the matched Kunitz comparison.

The manuscript fits \textit{JCIM} as a new molecular-modeling and chemical-informatics
methodology with explicit theory, reproducible computational benchmarks, and an honest
account of its domain of validity. It is not submitted as a docking or experimental binding
study. The low-confidence Cav2.2 complex predictions are retained in the Supporting
Information only as a negative diagnostic and are not used as evidence of binding.

The manuscript and Supporting Information provide a Data and Software Availability
statement. All curated inputs, canonical per-replicate outputs, benchmark emissions, scoring
results, and reproduction scripts are provided in the public study repository. The author
declares no competing financial interest.

A preprint is available as
\href{https://arxiv.org/abs/2603.20115}{arXiv:2603.20115}
(\href{https://doi.org/10.48550/arXiv.2603.20115}{doi:10.48550/arXiv.2603.20115}).
This manuscript is not under consideration elsewhere.

Thank you for your consideration.

\closing{Sincerely,}

\end{letter}

\end{document}
